\chapter*{Synthèse exécutive}
\addcontentsline{toc}{chapter}{Synthèse exécutive}

MABRID est conçu comme une place de marché numérique marocaine \emph{de confiance} : une plateforme multi-acteurs reliant acheteurs, vendeurs locaux et prestataires de services professionnels, tout en intégrant la vérification, la modération communautaire et une sécurité de niveau entreprise dans le modèle opérationnel. Cette synthèse exécutive condense l'ensemble du rapport à l'intention des décideurs, des investisseurs et des parties prenantes en matière de conformité.

\section*{Positionnement stratégique}

L'économie numérique marocaine se développe grâce à l'adoption du mobile, à la croissance du commerce électronique et aux programmes gouvernementaux soutenant l'entrepreneuriat. Toutefois, la fragmentation de la découverte des entreprises locales, le manque de signaux de confiance et l'inégale littératie numérique créent des frictions tant pour les consommateurs que pour les petites entreprises. MABRID comble cette lacune en combinant des fiches d'entreprises sélectionnées, des systèmes de réputation, des communautés par ville et une visibilité premium optionnelle pour les vendeurs---sans se positionner comme intermédiaire de paiement dans sa phase initiale.

\section*{Gouvernance et management}

L'entreprise fonctionne selon une structure de gouvernance définie séparant la supervision stratégique (conseil d'administration et direction exécutive), la gestion opérationnelle (produit, confiance et sécurité, finance) et les fonctions indépendantes de sécurité et de conformité. Les indicateurs clés de performance couvrent la liquidité de la place de marché (acheteurs et vendeurs actifs), les métriques de confiance (taux de vérification, délai de résolution des litiges), l'état de la sécurité (temps moyen de réponse aux incidents) et l'économie unitaire (coût d'acquisition client par rapport à la valeur vie client).

\section*{Architecture de sécurité}

La sécurité est traitée comme une priorité au niveau du conseil d'administration. MABRID adopte une stratégie \gls{zt} : accès fondé sur l'identité, moindre privilège, micro-segmentation et vérification continue. La gestion du cycle de vie des identités intègre le \gls{rbac}, le \gls{mfa} et des politiques d'accès conditionnel. Les données sont classifiées, chiffrées en transit et au repos, et sauvegardées avec des procédures de reprise testées. La gestion des menaces suit une évaluation structurée des risques, une modélisation des menaces, une gestion des vulnérabilités et un programme de réponse aux incidents intégré aux capacités \gls{siem} et \gls{soar} sur Azure (Microsoft Sentinel, Log Analytics).

\section*{Stratégie cloud et infrastructure}

Microsoft Azure constitue le socle cloud principal, gouverné par Azure Policy, la gestion des coûts, l'administration fondée sur les rôles et une discipline d'infrastructure en tant que code. Les objectifs de reprise après sinistre visent des seuils définis de temps de reprise et de point de reprise pour les services critiques de la place de marché. L'optimisation des coûts équilibre performance, résilience et viabilité financière lors du déploiement à l'échelle nationale.

\section*{Conformité juridique et éthique}

MABRID s'aligne sur la loi marocaine 09-08 relative à la protection des données à caractère personnel, les normes de protection des consommateurs et les exigences du commerce électronique, tout en appliquant les principes du \gls{gdpr} lorsque le traitement concerne des personnes concernées européennes. Les politiques de plateforme---conditions d'utilisation, règles vendeurs et acheteurs, modération de contenu et résolution des litiges---forment une couche de gouvernance intégrée. Les engagements éthiques comprennent la transparence, des pratiques équitables sur la place de marché, l'accessibilité et l'utilisation responsable de l'aide à la décision automatisée.

\section*{Feuille de route et perspectives}

Une feuille de route en trois phases progresse des villes d'ancrage vers une couverture nationale et une expansion internationale sélective. La viabilité financière repose sur des flux de revenus diversifiés et une maîtrise disciplinée des dépenses cloud. Les capacités futures comprennent une vérification d'identité renforcée, l'intelligence économique pour les vendeurs et des programmes de partenariat d'entreprise.

\section*{Recommandation}

Les parties prenantes devraient considérer ce rapport comme une documentation pré-lancement adaptée à la diligence raisonnable des investisseurs, aux discussions de partenariat d'entreprise et à l'engagement réglementaire auprès de la \gls{cndp}. Un examen juridique formel et une évaluation de sécurité indépendante demeurent des prérequis au lancement public.
